
%===========================================================
%    IN CLASS NOTES
%===========================================================
\subsection{Einstein's Equivalence Principle}
	\subsubsection{Inertial Frames}
		The Weinberg definition is, at every point in space time i.e $(t,x,y,z)$ in a gravitational field,
		it is possible to choose a locally inertial frame (i.e, free-falling)
		coordinate system such as the ``laws of nature'' take the same form
		as in an unaccelerated cartesian coordinate system in the absence of gravity.

		i.e., fields that are accelerating are are indistinguishable from gravity.
		
		That is, gravity fictitious forces, because it can be removed by a coordinate
		transformation. It is there only because we use the wrong coordinate system.

		The equivalence principle states that it is always possible to remove gravity,
		therefore gravity is a ``fictitious force" that can be ``removed" when writing
		the local field in any experiment.

		The difference in gravity and other fictitous forces is that gravity can only be
		removed localy, while other fictitious forces can be removed globally.
		We say that regardless of the mass of the object, it does not take part it its gravitational forces.

		The measurement of the differences between these gravitational fields is called tides.

		The usual way this is expressed is to imagine a person in an box and he conducts a free fall experiment,
		measuring some sort of experiment where he can measure accelerations. He will measure the acceleration of an object $i_1$
		\begin{equation}
			\vb{a}_1 = -\frac{m_1g}{m_1i} \vb{g}
		\end{equation}
		and for a second object $i_2$, he can measure
		\begin{equation}
			\vb{a}_2 = -\frac{m_1g}{m_1i} \vb{g}
		\end{equation}
		The universality of freefall is $a_1 = a_2$ is the same under any circumstances, called the universality of free fall.

		In that for, we can write the field as a non-inertial frame
		\begin{equation}
			\ddot{y} = -g
		\end{equation}

		We can remove gravity, however. Let $z$
		\begin{equation}
			z = y - \frac{1}{2}gt^2
		\end{equation}

		This will result to $\ddot z = 0$

		This is known as the inertial coordinate system, since there is no percieved acceleration. The point is,
		gravity can be removed locally.

		The person in the box cannot be able to tell if he is accelerating with constant accleration $g$ upwards,
		or under the influence of gravity. Gravity is simply the manifestation of being in the wrong frame.

		\textit{Inertial} means the absence of gravity.

		We have the weak equivalence principle, the laws of nature simply mean mechanics. For the strong equivalence principle, we mean that mechanics and
		strongly gravitating objects applies.

		Einstein equivalence principle, laws of nature means all laws of nature. It is the basis of general relativity.
		In any theory of gravity, you can remove $g$

		It means that gravity is genuinely fictitious if no experiment can tell that we are in a gravitational field.
		If one can remove gravity, then what remains should be the stage.

	\subsubsection{History of Relativity}
		The notion of inertiality plays a major role in relativity. Let us first explain relativity in its historical context.

		We begin with special relativity and Newtonian gravity. The combination of these two leads to gravity.

		\paragraph{Relativity} What is relativity? It is an old notion that exists even by the time of newton. There has to exists
		observers in which the laws of nature are the same. Regardless anywhere, the laws of nature is invariant of everywhere.
		This is called the equivalence of \textit{inertial observers}, and the laws of natuer are all the same to all inertial observers.

		Aristiotle did not have this notion. He believed that there is a special place where objects pull towards, the center of the universe, i.e,
		the center of the Earth. Anyone in this special region will get all the laws of nature. It makes sense to express physics in terms of that
		center of the Earth, since all frames are fixed relative to it.
		For example, waves in a pond are fixed relative to the center.

		For Galileo. In 1687 he rejected this notion by positing that there are observers for which the laws would be correct.
		This is called Galilean relativity. There are a class of observers which all move to relative each other at uniform velocities.

		\begin{align}
		 t' = t\\
		 \vb{x} &= \vb{x} - \vb{v}t
		\end{align}

		This relation explains that one can switch to the frame of reference of these two observers.
		In that case, one can have the following observation.
		\begin{equation}
			\vb{F} = m\vb{a} \implies \vb{F}' = m \vb{a}'
		\end{equation}

		Newtonian theory satisfied Galilean relativity, i.e., invariant under Galilean transformations.
		From 1687 onwards, the only fundamental laws of physics is Newtonian physics.

		For three centuries, this is the jewel of Physics. Then, electromagnetics broke it.

		In 1862, Maxwell arrived and showed that the Maxwell relations are not Galilean invariant, and such
		does not satisfy Galilean relativity, challenging hundreds of years of orthodoxy.
		They ended up with the conclusion that Maxwell's relations are not fundamental
		I.e., Maxwell's equations are only correct in the rest frame of something, and that thing is the aether,
		an invented medium where Maxwell's equation sits at. Maxwell's are waves in a pond.

		Light moves at speed $c$ only on the aether's rest frame. Thus, we have the Michelson-Morley experiment, which
		is a null experiment and the results say that $c$ is the same.
		People fought against this result (as it should, people should fight for theories to make the best
		theories right).

		We have Lorentz contraction -, unknown why it contracts, basta it contracts. adhoc solution to be consistent with the null result.

		Ether Dragging - imagien the ether not being a fixed frame, something like molasses or syrup, so the Earth in motion drags the molasses around it,
		and thus always remains locally at rest with its ether. Hence, there will be no changes in $c$.

		The debates went on for decades. People who built their carreers on it, so they made a lot of adhoc solutions.
		There are dozens of theories introduced in the hope of making Maxwell consistent with Galilean.

		Einstain came in 1905 with his annus mirabilis papers.

		The three papers gave bith to quantum theory, statistical mechanics, and special relativity. All of this before even
		turning 24.

		Everyone was so focused in meshing Galilean relativity and maxwells. The main idea is to deny the fundamentality of Maxwells
		in that it only works in the Eather.

		Einstein posits that Maxwell's is fundamental, and newtonian mechanics is not fundamental, a theory of three centuries old with so much backing.

		His solution is
		\begin{itemize}
			\item adopt $c$ as universal
			\item Accept Maxwell's equations are fundamental	
		\end{itemize}
		As such, the equations governing Galilean relativity must be dropped so that $c$ remains constant.
		By consequence, you also have to give up or change Newtonian mechanics and to drop it as fundamental.

		Divorce yourself from the notion that Newton is fundamental.

		The notion of universal observers are stull true, wehre there are observers where the laws of motion remain teh same, but to reconcile that and relativity,
		and to do that, one must unify space and time. Relativity and lorentz transformations.

		Maxwell is relativistic. Relativistically covariant, does not change its form as relativistic transformations is made
		Form invariant under the right transformations.

		task is to find the correct transformation laws that keep Maxwell invariant.

		Poincaire-Lorentz transformations were found to do that.
		\begin{align}
			ct'&= \gamma(ct-\beta x)\\
			x' &= \gamma(x - \beta c t)\\
			\gamma &= \frac{1}{\sqrt{1-\beta^2}} \\
			\beta &= \frac{v}{c}
		\end{align}

		The project, according to Einstein, is to rewrite all of physics (not too bad, the only physics then was Newtonian and Maxwells)
		so that it was in ``relativistically covariant" form. Invariant from one inertial observer to another.

		the form of these transformations is

		\begin{equation}
			x^{\mu'} = \Lambda^\mu_\nu x^\nu
		\end{equation}
		
		Where
		\begin{equation}
			x^\mu = (ct', x^1, x^2, x^3)
		\end{equation}

		and
		\begin{equation}
			\Lambda^\mu_ni = \bmat{
				\gamma & -\gamma \beta & 0 & 0\\
				-\gamma \beta & \gamma & 0 & 0 \\
				0 & 0 & 1 & 0 \\
				0 & 0 & 0 & 1\\
			}
		\end{equation}
		Now, the force equation becomes
		\begin{equation}
			\dv{\vb{p}}{t} = \vb{F}\implies \dv{p^\mu}{d\tau} = F^\mu
		\end{equation}
		The second equation limits to the first one in small values.

		\paragrph{How about gravity?}

		Graivtiy is just $\lpc \Phi = 4\pi G \rho$. This is not relativistically covariant, since it is only in space.

		Modify it so that
		\begin{equation}
			\square \Phi = 4\pi G T^\mu_{\ \mu}
		\end{equation}
		This is known as Nordstro gravity. But this is incorrect, since it does not
		predict the bending of light.

		However, Einstein was bothered with the equivalence principle. Why is there
		an equivalence of falling, i.e., gravity is peculiar.

		It was Minkowski in 1907 that placed these transformations in modern form.
		One can make sense of this by imagining space and time as a single structure, spacetime, with
		a metric nown known as Minkowski metric.
		\begin{equation}
			\eta_{\mu\nu} = \bmat{
				-1 & 0 & 0 & 0 \\
				0 & 1 & 0 & 0 \\
				0 & 0 & 1 & 0 \\
				0 & 0 & 0 & 1
			}
		\end{equation}
		and transformations that preserve the Minkowski metric is correct.

		In 1915, Einstein formalized his theory of General Relativity, where he formalized
		the notion of space time, and space time need not only be the Minkowski metric.
		\begin{equation}
			\eta_{\mu\nu}\to	g_{\mu\nu}
		\end{equation}
		He expanded the notion of inertial observers to all observers in free fall
		under these space time, implying that inertial observers are no longer just
		moving with uniform velocity to each other, but all observers that are freely
		falling, i.e., geodesics.

		He got it alll quite early, but it was getting the Field Equations that took
		a long time.

		Relativity is very very old, and the only thing that Einstein did is expand it to an absolute structure.
		Minkowski's Relativity is special relativity $\eta$, called special because it is
		restricted to a class of observers moving constantly relative to each other,
		and the general form $g$ became general relativity, with gravity.

		\textit{machian}, gr is not \textit{machian}, everything is relational, nothing should be absolute.

		There is a notion that light is inherently tide to relativity. However, if Maxwell's break, relativity
		remains true. Light has nothing to do with relativity, it is a historical bearing, but it has no bearing to relativity. Relativity is a metaprinciple. Maxwell inspiried it, but is conceptually distinct. W
		Light is not tied to relativity, and the next lecture goes on to deriving GR that makes no mention
		of light but gets the same equations.

		The only role of Maxwell is to define $c$ to be a number.
		We can approach a lot of relativity to different numbers. If we assign $c= \infty$, you get Galilean relativity,
		if you assign $c=0$, you get
		Carolian relativity. And if you assign $c$ to be an arbitrary number between $[0,\infty)$, you get General Relativity.
